\subsection{Các mô hình sharing economy}
\subsubsection*{GrabBike / BeBike (dịch vụ gọi xe 2 bánh)}
\textbf{Điểm mạnh:}
\begin{itemize}
    \item Nền tảng lớn, UX/UI quen thuộc.
    \item Tích hợp thanh toán điện tử, bản đồ định vị chính xác.
\end{itemize}

\textbf{Điểm yếu:}
\begin{itemize}
    \item Không phải P2P (không chia sẻ tài sản cá nhân, mà là dịch vụ vận tải).
    \item Chi phí vận hành cao (phần trăm chiết khấu).
    \item Người dùng không được tự do chọn xe, chỉ chọn chuyến.
\end{itemize}

\subsubsection*{Xanh SM (VinFast – dịch vụ taxi/xe điện)}
\textbf{Điểm mạnh:}
\begin{itemize}
    \item Hệ sinh thái đồng bộ (xe, trạm sạc, quản lý vận hành).
    \item Hình ảnh thương hiệu mạnh, tạo niềm tin.
\end{itemize}

\textbf{Điểm yếu:}
\begin{itemize}
    \item Mô hình tập trung, cần đầu tư hạ tầng lớn.
    \item Không hướng tới cộng đồng nhỏ lẻ (ví dụ chung cư, ký túc xá).
\end{itemize}

\subsubsection*{Airbnb (mô hình P2P sharing trong lĩnh vực lưu trú)}
\textbf{Điểm mạnh:}
\begin{itemize}
    \item Thành công toàn cầu trong sharing economy.
    \item Cơ chế đánh giá 2 chiều $\rightarrow$ xây dựng trust.
\end{itemize}

\textbf{Điểm yếu:}
\begin{itemize}
    \item Chưa áp dụng cho phương tiện ở Việt Nam.
    \item Phụ thuộc rất nhiều vào quản lý cộng đồng \& chính sách địa phương.
\end{itemize}

\subsubsection*{Hệ thống xe đạp điện công cộng (Hà Nội, TP.HCM)}
\textbf{Điểm mạnh:}
\begin{itemize}
    \item Quản lý tập trung, dễ giám sát.
    \item Chi phí rẻ cho người dùng.
\end{itemize}

\textbf{Điểm yếu:}
\begin{itemize}
    \item Hạn chế khu vực, không linh hoạt.
    \item Không tận dụng tài sản nhàn rỗi của cư dân.
\end{itemize}

\begin{longtable}{|p{2.2cm}|p{3.5cm}|c|c|c|c|}
\caption{So sánh các mô hình tiêu biểu trong lĩnh vực chia sẻ và di chuyển} \\
\hline
\textbf{Mô hình} & \textbf{Loại hình} & \textbf{Tính P2P} & \textbf{Chi phí} & \textbf{Linh hoạt} & \textbf{Mở rộng} \\ \hline
\endfirsthead

\multicolumn{6}{c}%
{\tablename\ \thetable{} -- \textit{Tiếp theo}} \\
\hline
\textbf{Mô hình} & \textbf{Loại hình} & \textbf{Tính P2P} & \textbf{Chi phí} & \textbf{Linh hoạt} & \textbf{Mở rộng} \\ \hline
\endhead

\hline \multicolumn{6}{r}{\textit{(Còn tiếp)}} \\
\endfoot

\hline
\endlastfoot

GrabBike / BeBike & Dịch vụ gọi xe 2 bánh & ⭐☆☆☆☆ & ⭐⭐☆☆☆ & ⭐⭐☆☆☆ & ⭐⭐⭐⭐☆ \\ \hline
Xanh SM (VinFast) & Taxi/xe điện tập trung & ☆☆☆☆☆ & ⭐⭐⭐⭐⭐ & ⭐☆☆☆☆ & ⭐⭐☆☆☆ \\ \hline
Airbnb & Chia sẻ lưu trú & ⭐⭐⭐⭐⭐ & ⭐⭐☆☆☆ & ⭐⭐⭐⭐⭐ & ⭐⭐⭐⭐⭐ \\ \hline
Xe đạp điện công cộng & Dịch vụ công cộng & ⭐☆☆☆☆ & ⭐⭐☆☆☆ & ⭐☆☆☆☆ & ⭐⭐☆☆☆ \\ \hline
\textbf{Đề tài nghiên cứu} & \textbf{Chia sẻ xe máy điện (P2P)} & \textbf{⭐⭐⭐⭐☆} & \textbf{⭐☆☆☆☆} & \textbf{⭐⭐⭐⭐☆} & \textbf{⭐⭐⭐⭐☆} \\ \hline

\end{longtable}

\textbf{Kết luận:} Chưa có hệ thống nào tại Việt Nam kết hợp \textbf{P2P model} cho xe máy điện. Đây chính là khoảng trống thị trường và cơ hội nghiên cứu/triển khai.
